\newcommand{\ToneQualityTxCostCaption}{Quality and online transmitter software time at $N=1024$, separately at4,10 and19~dB.
The horizontal axis reports the arithmetic mean (primary panels) or median (companion panels) of actual measured TX time in milliseconds on a logarithmic scale.
TX spans the input image tensor to the full I/Q frame, including visual encoding, necessary probability computation, entropy encoding and length-based fallback attempts, FEC and modulation; model loading, file I/O, metrics, queueing and channel propagation are excluded.
The same NVIDIA GeForce RTX4090D environment, batch size1, frozen numerical settings and GPU synchronization protocol are used for all four systems. Each point aggregates16 fixed development sources with three measured repetitions after the registered warm-ups (48 measured calls); these small-sample timings are descriptive, with no timing confidence interval or universal deployment/tail-latency claim.
Quality means and vertical error bars instead come from the previously examined500-source holdout, three noises per source, with existing pointwise95\% source-bootstrap intervals and no multiple-comparison correction.
This is a post-hoc comparison of method/workpoint aggregates across two different populations, not an image-paired quality-versus-cost correlation. All scheduled failure outcomes remain included. ConvNeXt reports prediction agreement, not accuracy.
The original raw complete-scale and partial-scale policies and intervals are unchanged; the complete-scale entropy policies were selected on calibration data. The plotted protocols retain their distinct raw KEEP versus entropy rejection/fallback rules. No quality value, confidence interval or timing is synthesized for an unmeasured setting.}
